%% Chapter 21 ----------------------------------------------------------------
\chapter{Fitting the Flux Rope and the Shock}
\label{ch:gcs-fitting}

This chapter describes the models of the GCS CME Fitting window, how they are aligned and
refined against the observed fronts, how fits are recorded over time to build kinematics, and
how results are exported.

\section{The Model card}

The \gui{Model} card (Figure~\ref{fig:gcs-model-card}) shows the parameters of one model at a
time. Choose the model with \gui{Edit} on the card's title row:
\begin{description}
  \item[GCS flux rope] the Graduated Cylindrical Shell, fitted to the CME ejecta.
  \item[Shock] a spheroid or ellipsoid, fitted to the shock the CME drives --- its faint outer
    envelope (Section~\ref{sec:shock-fitting}).
\end{description}
Both models are drawn in every panel (the shell in orange and the shock in sky blue by
default); the edited model is the one that the handles, clicks, \gui{Refine fit}\index{GCS!refine fit},
\gui{Commit} and the Kinematics card act on.

\screenshot[0.45\linewidth]{gcs/gcs_model_card}{The \gui{Model} card with the GCS parameters.}{fig:gcs-model-card}{The Model card in GCS flux rope mode: the six sliders with numeric entry (Lon, Lat, Tilt, Height, α, κ), the derived-geometry readout, Refine fit and Commit GCS, and the Pick front points, wireframe, Undo point and Clear points controls.}

\section{GCS parameters}
\label{sec:gcs-parameters}
\index{GCS!parameters}

Six sliders, each with a numeric entry, control the shell (Table~\ref{tab:gcs-params}).

\begin{table}[!htbp]
  \centering
  \caption{Parameters of the GCS flux-rope model.}
  \label{tab:gcs-params}
  \small
  \begin{tabularx}{\linewidth}{@{}P{0.12\linewidth}P{0.2\linewidth}L@{}}
    \toprule
    \thead{Slider} & \thead{Range} & \thead{Meaning} \\
    \midrule
    Lon & \(-180^\circ\) to \(180^\circ\) & Stonyhurst longitude of the direction of propagation. \\
    Lat & \(-89^\circ\) to \(89^\circ\) & Stonyhurst latitude of the direction of propagation. \\
    Tilt & \(-90^\circ\) to \(90^\circ\) & Rotation of the shell about the propagation direction. \\
    Height & 1.05--30~\(R_\odot\) & Distance of the leading edge (apex) from Sun center --- not the altitude above the surface. \\
    \(\alpha\) & \(1^\circ\)--\(80^\circ\) & Half-angle between the legs of the shell. \\
    \(\kappa\) & 0.05--0.95 & Aspect ratio, controlling the thickness of the shell. \\
    \bottomrule
  \end{tabularx}
\end{table}

Set the direction, tilt, \(\alpha\) and \(\kappa\) before fine-tuning the height. The apex can
also be dragged in any panel: moving it around the Sun turns the direction, and moving it
outwards or inwards changes the height. The readout below the sliders gives the derived leg
height, apex cross-section radius, center distance, and the intrinsic face-on and edge-on
widths.

\begin{background}[Background: GCS geometry]
  The GCS is a hollow croissant: two conical legs joined by a curved front with a circular
  cross-section that grows with distance from the Sun \parencite{thernisien2006,thernisien2011}.
  The face-on width is \(2[\alpha + \arcsin\kappa]\) and the edge-on width \(2\arcsin\kappa\);
  both include the shell thickness and differ from the apparent width in a coronagraph image.
  Because the height is the leading edge, it feeds the height--time fit directly. GCS is a
  flux-rope model, not a shock model.
\end{background}

\section{Clicking the front}
\label{sec:front-points}

With \gui{Pick front points}\index{GCS!front points} ticked, left-click points along the same front in two or more
views. Points are drawn in each panel and are used by \gui{Refine fit}:
\begin{itemize}
  \item a right-click removes that panel's last point;
  \item \gui{Undo point} (\kbd{Ctrl+Z}) removes the point clicked last in any panel;
  \item \gui{Clear points} removes all of the edited model's points on the displayed images;
  \item untick \gui{Pick front points} to navigate without adding points.
\end{itemize}
The \gui{GCS wireframe} and \gui{Shock wireframe} options hide a model temporarily so that the
observed front can be inspected. Each model keeps its own front points.

\section{Refining the fit}
\label{sec:refine}
\index{GCS!refinement}

\gui{Refine fit} (\kbd{Ctrl+R}) makes a local least-squares adjustment of the edited model
through the clicked points. It refines a fit; it cannot find one from scratch, because
longitude, \(\alpha\) and \(\kappa\) are poorly conditioned in GCS, so first align the shell
roughly with the CME by eye. The refinement fits as many parameters as the points support
(Table~\ref{tab:refine-stages}); the status bar names the parameters fitted and how many more
points would free the next one.

\begin{table}[!htbp]
  \centering
  \caption{Staged refinement of the GCS model.}
  \label{tab:refine-stages}
  \small
  \begin{tabularx}{\linewidth}{@{}P{0.38\linewidth}L@{}}
    \toprule
    \thead{Front points} & \thead{Parameters fitted} \\
    \midrule
    2 & Height \\
    4, in two separated views & Height, longitude, latitude \\
    5 & \ldots and tilt \\
    6 & \ldots and \(\alpha\) \\
    7 & All six parameters (adding \(\kappa\)) \\
    \bottomrule
  \end{tabularx}
\end{table}

The status bar also reports the views taking part, their angular separations and the number
of points. Only images within the \gui{Max time offset}\index{GCS!max time offset} of the shared time take part. With a
single view, or with no pair of views between 20° and 160° apart, the direction cannot be
constrained and is held fixed; empty, nearly coincident and nearly opposite viewpoints do not
constrain the direction independently. Formal errors are shown when the local solution is
reliable, are withheld otherwise, and become invalid when the model or the observations
change.

\section{Recording fits}
\label{sec:gcs-recording}

\gui{Commit GCS}\index{GCS!commit} (\kbd{Ctrl+Return}) records the current fit at the shared time, together with
the images it was made from and their observation details. Step to later times and repeat to
build a height--time series. One combination of images cannot be recorded at two different
times, so repeated images never count as independent height measurements.

When stepping, dragging the time slider or playing, the recorded shells\index{GCS!recorded shells} are drawn: each fit on
its own images, interpolated between recorded times (a display, not a fit) and held before the
first and after the last record. A model with nothing recorded keeps its slider values, so
commit before stepping away. \menu{Fit > Restore recorded model at current time} and
\menu{Delete recorded fit at current time} revisit or remove a record; \menu{Fit > Reset model
parameters} resets the edited model.

Moving the event range to a different event sets the recorded fits aside. They return when you
come back to that event, and the JSON export includes them.

\section{Kinematics}
\label{sec:gcs-kinematics}
\index{GCS!kinematics}

The \gui{Kinematics} card (Figure~\ref{fig:gcs-kinematics}) lists the recorded fits of the
edited model beside a height--time graph. For the GCS model the table shows \gui{Time (UT)},
\gui{t (s)}, \gui{Apex (R☉)}, \gui{Lon (°)}, \gui{Lat (°)} and \gui{α (°)}; for the shock,
\(\kappa\) replaces \(\alpha\).

\screenshot{gcs/gcs_kinematics}{A GCS fit with three recorded times; the \gui{Kinematics} card (right) lists the recorded fits beside their height--time graph and linear fit.}{fig:gcs-kinematics}{The Kinematics · GCS card with several recorded fits in the table, the height–time graph with a quadratic fit and error bars, and the Fit, Fit height–time, Clear fits and Export controls.}

\begin{itemize}
  \item Double-click a row to go back to its time and restore its model.
  \item Choose a \gui{Linear}\index{height--time fit!order}, \gui{Quadratic} or \gui{Cubic} fit and click
    \gui{Fit height–time}\index{GCS!kinematics}. A fit of degree \(n\) needs \(n+1\) recorded times, and \(n+2\)
    before errors can be estimated (Section~\ref{sec:tracking-panel}).
  \item \gui{Clear fits} removes the recorded fits of the edited model.
  \item \gui{Export} saves the table as CSV or the graph as a figure.
\end{itemize}
The heights are de-projected under the model assumptions, so the speeds are model-dependent
radial speeds rather than plane-of-sky speeds.

\section{Fitting the shock}
\label{sec:shock-fitting}
\index{shock model}\index{spheroid}\index{ellipsoid}

Choose \gui{Edit: Shock} to fit the shock with a \gui{Spheroid} or an \gui{Ellipsoid}
(Figure~\ref{fig:gcs-shock}). The parameters follow PyThea\index{PyThea}'s conventions exactly
(Table~\ref{tab:shock-model-params}).

\begin{table}[!htbp]
  \centering
  \caption{Parameters of the shock model.}
  \label{tab:shock-model-params}
  \small
  \begin{tabularx}{\linewidth}{@{}P{0.13\linewidth}P{0.2\linewidth}L@{}}
    \toprule
    \thead{Slider} & \thead{Range} & \thead{Meaning} \\
    \midrule
    Lon, Lat & as for GCS & Stonyhurst direction of the shock center and apex. \\
    Height & 1.05--30~\(R_\odot\) & Heliocentric distance of the apex (center distance plus the radial semi-axis \(a\)). \\
    \(\kappa\) & 0.05--2.0 & Self-similar constant \(\kappa = b/(h - 1\,R_\odot)\): the lateral size for a given apex height. \\
    \(\varepsilon\) & \(-0.99\) to \(0.99\) & Signed eccentricity: positive stretches the shock radially (\(a > b\)), negative flattens it (\(a < b\)). \\
    Tilt & \(-90^\circ\) to \(90^\circ\) & Ellipsoid only: rotation about the radial axis; at 0 the \(b\) axis is parallel to the solar equator. \\
    \(\alpha\) (b/c) & 0.5--1.5 & Ellipsoid only: ratio of the two lateral semi-axes. \\
    \bottomrule
  \end{tabularx}
\end{table}

\screenshot{gcs/gcs_shock_fit}{Shock and flux-rope models in three viewpoints.}{fig:gcs-shock}{The three viewpoints with both the GCS shell (orange) and the shock spheroid (sky blue) drawn, Edit: Shock selected, and shock front points clicked on the outer envelope.}

The derived readout gives the center distance and the semi-axes, under PyThea's names. Click the
shock front and use \gui{Refine fit}: the model's exactly computed projected outline is pulled
through the points in every view, with its own sequence of height, direction and shape
parameters. Record shock fits with \gui{Commit Shock}\index{shock model!commit}. The shock keeps its own front points,
recorded fits, kinematics and CSV file. With fewer than three views, \(\varepsilon\) and the tilt
are weakly constrained and an ellipsoid's shape parameters trade off against each other, so
prefer a spheroid.

\section{Exporting results}
\label{sec:gcs-export}
\index{GCS!export}

Images and graphs are drawn again from the data rather than copied from the screen, always on
a white page whatever the application theme; graphs follow the OriginPro style. Figures can be
saved as PNG, PDF, EPS, SVG, TIFF or JPG.
\begin{description}
  \item[Export analysis JSON…] (\kbd{Ctrl+Shift+S}) the current, recorded and set-aside
    parameters of both models, with the clicked points, image times, offsets and observer
    geometry. The file (schema version~2) has a \texttt{shock} section beside the GCS keys. It
    is an analysis export, not a PyThea session file.
  \item[Export recorded fits CSV…] the recorded series of the edited model
    (\file{cme_gcs_fit.csv} or \file{cme_shock_fit.csv}).
  \item[Export height–time graph…] the recorded apex heights with error bars and the chosen
    fit, titled with the model and fit, with the speed and acceleration in the legend.
  \item[Save viewpoint snapshot…] the three views as shown, with arcsecond axes, the
    wireframes, the solar limb, the front points and a note of where each drawn shell comes
    from. With the banners hidden, only each panel's name is printed above it.
  \item[Export movie (GIF/MP4)…] every time step as playback shows it, recorded shells
    included, at the playback speed. MP4 needs the bundled FFmpeg; a GIF is offered when it is
    unavailable.
  \item[Export fitting report (PDF)…] the viewpoints and their separations, the model at the
    current time, every recorded fit drawn on the images it was made from with its parameters
    and observation details, each model's linear, quadratic and cubic height--time fits (each
    with a titled graph and all its parameters, then compared), and the method, its limits and
    references (Figure~\ref{fig:gcs-report}).
\end{description}
Stepping through time for a movie or a report leaves the window as it was: the time shown, the
working models, their formal errors and the front points are restored afterwards.

\screenshot[0.55\linewidth]{gcs/gcs_report_page}{A page of the GCS fitting report.}{fig:gcs-report}{A representative page of the exported PDF fitting report showing a recorded fit drawn on its three images with the parameter table.}

\begin{caution}[Scientific limits]
  GCS models the flux rope, not the shock. Formal fit errors leave out the uncertainty from
  identifying the front, non-simultaneous images, image preparation and the assumed geometry,
  and a small residual does not make a fit unique or accurate \parencite{verbeke2023}. Helioviewer\index{Helioviewer}
  JPEG2000 images are display products: suitable for fitting shapes, not for calibrated
  intensity measurements. \menu{Help > Fitting workflow and scientific limits…} (\kbd{F1})
  summarizes the procedure and these limits.
\end{caution}
